% الجزء: sets-functions-relations
% الفصل: relations
% القسم: relations-as-sets

\documentclass[../../../include/open-logic-section]{subfiles}

\begin{document}

\olfileid[ar]{sfr}{rel}{set}
\olsection{العلاقات بوصفها مجموعات}

\begin{explain}
ذكرنا في \olref[sfr][set][imp]{sec} بعض المجموعات المهمة:
$\Nat$ و$\Int$ و$\Rat$ و$\Real$. ولا ريب أنك تتذكر بعض العلاقات
المثيرة للاهتمام بين !!{element}s بعض هذه المجموعات. فلكل واحدة من هذه
المجموعات، مثلًا، \emph{علاقة ترتيب} معيارية تمامًا. وهناك أيضًا علاقة
\emph{«مطابق لـ»}، التي تقوم بين كل كائن ونفسه ولا تقوم بينه وبين أي شيء
آخر. وسنصادف علاقات أخرى كثيرة مثيرة للاهتمام، بل إن العلاقات الممكنة
أكثر من ذلك. وقبل أن نستعرضها، سنبدأ بالإشارة إلى أنه يمكن النظر إلى
العلاقات بوصفها نوعًا خاصًا من المجموعات.

ولهذا الغرض، تذكّر أمرين من \olref[sfr][set][pai]{sec}. أولًا، تذكّر
مفهوم \emph{الزوج المرتب}: فإذا أعطينا $a$ و$b$، أمكننا تكوين
~$\tuple{a, b}$. ومن المهم أن ترتيب العناصر \emph{معتبر} هنا. ولذلك، إذا
كان $a \neq b$، فإن $\tuple{a, b} \neq \tuple{b, a}$. (قارن ذلك
بالأزواج غير المرتبة، أي المجموعات التي عدد عناصرها $2$، حيث
$\{a, b\}=\{b, a\}$.) ثانيًا، تذكّر مفهوم \emph{الضرب الديكارتي}: فإذا
كانت $A$ و$B$ مجموعتين، أمكننا تكوين~$A \times B$، وهي مجموعة جميع
الأزواج $\tuple{x, y}$ التي فيها $x \in A$ و$y \in B$. وبوجه خاص،
$A^{2}= A \times A$ هي مجموعة جميع الأزواج المرتبة من~$A$.

سننظر الآن في علاقة بعينها على مجموعة: العلاقة $<$ على مجموعة~$\Nat$
للأعداد الطبيعية. اعتبر مجموعة جميع أزواج الأعداد $\tuple{n, m}$ التي
فيها $n<m$، أي:
\[
  R=\Setabs{\tuple{n, m}}{n, m \in \Nat \text{ و } n<m}.
\]
هناك ارتباط وثيق بين كون $n$ أصغر من $m$ وكون الزوج $\tuple{n, m}$
عنصرًا في~$R$، وهو:
\[
      n<m\text{ إذا وفقط إذا }\tuple{n, m} \in R.
\]
والواقع أنه يمكننا، من دون أي فقد للمعلومات، أن نعد المجموعة $R$
\emph{هي} العلاقة $<$ على~$\Nat$.

وبالطريقة نفسها، يمكننا أن ننشئ لكل علاقة بين الأعداد مجموعة جزئية من
$\Nat^{2}$. وبالعكس، إذا أعطيت أي مجموعة من أزواج الأعداد
$S \subseteq \Nat^{2}$، كانت هناك علاقة مناظرة بين الأعداد، هي العلاقة
التي تقوم بين $n$ و$m$ إذا وفقط إذا كان $\tuple{n, m} \in S$. وهذا يسوغ
التعريف الآتي:
\end{explain}

\begin{defn}[العلاقة الثنائية]
\emph{العلاقة الثنائية} على مجموعة $A$ هي مجموعة جزئية من~$A^{2}$. فإذا
كانت $R \subseteq A^{2}$ علاقة ثنائية على~$A$ وكان $x, y \in A$، كتبنا
أحيانًا $Rxy$ (أو $xRy$) بدلًا من $\tuple{x, y} \in R$.
\end{defn}

\begin{ex}
  \ollabel{relations}
يمكن سرد مجموعة أزواج الأعداد الطبيعية $\Nat^{2}$ في مصفوفة ثنائية الأبعاد
كما يأتي:
\[
  \begin{array}{ccccc}
  \mathbf{\tuple{ 0,0 }} & \tuple{ 0,1 } &
    \tuple{ 0,2 } & \tuple{ 0,3 } & \ldots\\
  \tuple{ 1,0 } & \mathbf{\tuple{ 1,1 }} &
    \tuple{ 1,2 } & \tuple{ 1,3 } & \ldots\\
  \tuple{ 2,0 } & \tuple{ 2,1 } &
    \mathbf{\tuple{ 2,2 }} & \tuple{ 2,3 } & \ldots\\
  \tuple{ 3,0 } & \tuple{ 3,1 } & \tuple{ 3,2 } &
    \mathbf{\tuple{ 3,3 }} & \ldots\\
  \vdots & \vdots & \vdots & \vdots & \mathbf{\ddots}
  \end{array}
\]
كتبنا القطر هنا بخط عريض، لأن المجموعة الجزئية من $\Nat^2$ المكونة من
الأزواج الواقعة على القطر، أي:
\[
  \{\tuple{0,0 }, \tuple{ 1,1 }, \tuple{ 2,2 }, \dots\},
  \]
هي \emph{علاقة الهوية على}~$\Nat$. (ولشيوع علاقة الهوية، نعرّف
$\Id{A}=\Setabs{\tuple{ x,x }}{x \in A}$ لأي مجموعة $A$.) أما المجموعة
الجزئية المكونة من جميع الأزواج الواقعة فوق القطر، أي:
\[
  L = \{\tuple{ 0,1 },\tuple{ 0,2 },\ldots,\tuple{ 1,2 },
  \tuple{ 1,3 }, \dots, \tuple{ 2,3 }, \tuple{ 2,4 },\ldots\},
\]
فهي علاقة \emph{أصغر من}، أي إن $Lnm$ إذا وفقط إذا كان $n<m$. والمجموعة
الجزئية المكونة من الأزواج الواقعة تحت القطر، أي:
\[
  G=\{ \tuple{ 1,0 },\tuple{ 2,0 },\tuple{
    2,1 }, \tuple{ 3,0 },\tuple{ 3,1 },\tuple{ 3,2 }, \dots\},
\]
هي علاقة \emph{أكبر من}، أي إن $Gnm$ إذا وفقط إذا كان $n>m$. واتحاد
$L$ مع $I$، الذي يمكن أن نسميه $K=L\cup I$، هو علاقة \emph{أصغر من أو
يساوي}: إذ يكون $Knm$ إذا وفقط إذا كان $n \le m$. وبالمثل، فإن
$H=G \cup I$ هي \emph{علاقة أكبر من أو يساوي.} والعلاقات $L$ و$G$ و$K$
و$H$ أنواع خاصة من العلاقات تسمى \emph{ترتيبات}. ولـ$L$ و$G$ الخاصية
الآتية: لا يرتبط أي عدد بنفسه بالعلاقة $L$ أو بالعلاقة $G$ (أي إنه، لكل
$n$، لا $Lnn$ ولا $Gnn$). وتسمى العلاقات ذات هذه الخاصية
\emph{علاقات لا انعكاسية}؛ فإذا
كانت أيضًا ترتيبات، سميت \emph{ترتيبات صارمة.}
\end{ex}

\begin{explain}
مع أن الترتيبات وعلاقة الهوية علاقات مهمة وطبيعية، ينبغي التشديد على أن
تعريفنا يعد \emph{أي} مجموعة جزئية من $A^{2}$ علاقة على~$A$، مهما بدت
غير طبيعية أو مصطنعة. وبوجه خاص، فإن $\emptyset$ علاقة على أي مجموعة
(وهي \emph{العلاقة الخالية} التي لا تقوم بين أي زوج من العناصر)، كما أن
$A^{2}$ نفسها علاقة على~$A$ أيضًا (وهي علاقة تقوم بين كل زوج)، وتسمى
\emph{العلاقة الشاملة}. بل إن مجموعة مثل
$E=\Setabs{\tuple{n, m}}{n>5 \text{ أو } m \times n \ge 34}$ تعد علاقة
أيضًا.
\end{explain}

\begin{prob}
  اسرد !!{element}s العلاقة $\subseteq$ على المجموعة
  $\Pow{\{a, b, c\}}$.
\end{prob}

\end{document}
